\phant{PHƯƠNG PHÁP ĐÁNH GIÁ}\label{sec:phuongphap}

\subsection{Thiết kế tổng thể}

Đơn vị phân tích là một cặp (mô hình ngôn ngữ, biến thể lập trình). Với bốn mô
hình và năm biến thể, thiết kế có $4\times5=20$ chương trình. Mỗi chương trình
được sinh từ một prompt cố định, biên dịch, chạy, sửa lỗi tối thiểu và đo đạc
theo cùng một quy trình. Cách bố trí này cho phép so sánh theo hai chiều: giữa
các mô hình trên cùng một biến thể, và giữa các biến thể trong cùng một mô hình.

\begin{table}[H]
\centering
\small
\caption{Ma trận thực nghiệm: bốn mô hình $\times$ năm biến thể.}
\label{tab:matrix}
\begin{tabularx}{\textwidth}{@{}l *{4}{>{\centering\arraybackslash}X}@{}}
\toprule
\textbf{Biến thể} & \textbf{ChatGPT-4} & \textbf{ChatGPT-5} & \textbf{Claude} & \textbf{LLaMA} \\
\midrule
C++11 bất đồng bộ        & C11 & C11 & C11 & C11 \\
C++17 thuật toán song song & C17 & C17 & C17 & C17 \\
C++20 coroutine          & C20 & C20 & C20 & C20 \\
OpenMP                   & OMP & OMP & OMP & OMP \\
MPI+OpenMP               & MPI & MPI & MPI & MPI \\
\bottomrule
\end{tabularx}
\end{table}

Mỗi ô là một lần sinh độc lập. Bài báo không lặp lại nhiều lần cho mỗi ô, không
thay đổi nhiệt độ hay cách diễn đạt prompt. Hệ quả là kết quả mô tả \emph{một}
mẫu đầu ra tiêu biểu, không phải phân phối đầy đủ của mô hình. Đổi lại, thiết kế
giữ được tính nhất quán: mọi so sánh đều dựa trên cùng một điều kiện sinh mã.

\subsection{Yêu cầu chức năng và tiêu chí chấp nhận}

Mỗi chương trình phải thỏa mãn đồng thời các yêu cầu sau:

\begin{enumerate}[label=\textbf{Y\arabic*.},leftmargin=1.6cm]
  \item Biên dịch được bằng GCC 11.5.0 trên nền Linux của cụm thử nghiệm.
  \item Chạy được với tham số dòng lệnh \code{-t <số luồng>}; riêng bản MPI dùng
        \code{-p <số phân vùng>}.
  \item Xuất ảnh Mandelbrot ra tệp theo định dạng PBM/PGM, không cần thư viện
        đồ họa ngoài.
  \item In thời gian chạy \emph{không bao gồm} I/O.
  \item Cho kết quả hình ảnh đúng: tập Mandelbrot hiện rõ, không đen toàn bộ,
        không lệch màu do sai định dạng.
\end{enumerate}

Ba tiêu chí đánh giá được định nghĩa như sau:

\begin{itemize}
  \item \textbf{Biên dịch (build)}: chương trình biên dịch thành công mà không
        cần sửa mã nguồn.
  \item \textbf{Thực thi (runtime)}: chương trình chạy đến khi kết thúc mà không
        treo, không crash, không báo lỗi bộ nhớ.
  \item \textbf{Tính đúng (correctness)}: ảnh đầu ra đúng về mặt trực quan so với
        ảnh tham chiếu.
\end{itemize}

Cần nhấn mạnh rằng “biên dịch được” không kéo theo “chạy đúng”, và “chạy được”
không kéo theo “tôn trọng số luồng”. Ba tiêu chí này độc lập với nhau; khoảng
cách giữa chúng chính là một phần quan trọng của kết quả.

\subsection{Thiết kế prompt}

Ba prompt được dùng nguyên văn cho mọi mô hình. Prompt P1 dùng chung cho ba biến
thể bộ nhớ chia sẻ, chỉ thay phần trong ngoặc nhọn; P2 dùng cho OpenMP; P3 dùng
cho MPI+OpenMP.

\begin{lstlisting}[language={},caption={Prompt P1 -- bộ nhớ chia sẻ với C++11/17/20.},label={lst:p1},numbers=none,basicstyle=\ttfamily\small]
Write a parallel code for the Mandelbrot set using C++ {11 async future,
17 parallel algorithms, 20 coroutines} and write the output to the PBM
format. Make the number of threads configurable with the command line
option -t and print the runtime excluding IO.
\end{lstlisting}

\begin{lstlisting}[language={},caption={Prompt P2 -- OpenMP.},label={lst:p2},numbers=none,basicstyle=\ttfamily\small]
Write a parallel code for the Mandelbrot set using C++ using OpenMP and
write the output to the PBM format. Make the number of threads
configurable with the command line option -t and print the runtime
excluding IO.
\end{lstlisting}

\begin{lstlisting}[language={},caption={Prompt P3 -- MPI+OpenMP.},label={lst:p3},numbers=none,basicstyle=\ttfamily\small]
Write a distributed code using MPI for the Mandelbrot set and using
OpenMP for parallelism in C++ and write the output to the PBM format.
Make the number of partitions configurable with the command line
option -p and print the runtime excluding IO.
\end{lstlisting}

Mỗi mệnh đề trong prompt có một mục đích kiểm thử:

\begin{itemize}
  \item \textbf{“using C++ \{11, 17, 20\}”}: buộc mô hình chọn đúng cơ chế của
        từng phiên bản chuẩn, thay vì dùng OpenMP cho mọi trường hợp.
  \item \textbf{“PBM format”}: tạo đầu ra có thể kiểm tra bằng mắt và bằng tệp,
        đồng thời loại bỏ phụ thuộc vào thư viện đồ họa.
  \item \textbf{“-t” / “-p”}: kiểm tra khả năng phân tích dòng lệnh và — quan
        trọng hơn — khả năng \emph{sử dụng} tham số đó trong phần tính toán.
  \item \textbf{“excluding IO”}: tách chi phí tính toán khỏi chi phí ghi ảnh, cần
        thiết cho phép đo mở rộng.
\end{itemize}

Với MPI, prompt ban đầu chưa đủ. Nếu không nêu rõ “số phân vùng cấu hình được”,
mã sinh ra có xu hướng cho mọi tiến trình tính lại toàn bộ tập Mandelbrot rồi
chỉ ghi phần của mình. Chương trình khi đó trông như song song nhưng không chia
việc, và hiệu năng không tăng theo số nút. Việc bổ sung ràng buộc này là một ví
dụ cụ thể cho thấy chất lượng đặc tả ảnh hưởng trực tiếp đến chất lượng mã.

\subsection{Quy trình kiểm định}

Quy trình gồm năm bước, áp dụng giống nhau cho mọi chương trình:

\begin{enumerate}
  \item \textbf{Biên dịch}: dùng GCC 11.5.0, ghi nhận toàn bộ lỗi và cảnh báo.
  \item \textbf{Sửa lỗi biên dịch}: tác giả sửa với thay đổi tối thiểu, chỉ để
        chương trình biên dịch được; không thiết kế lại thuật toán.
  \item \textbf{Thực thi}: chạy với một luồng và nhiều luồng (hoặc một và nhiều
        tiến trình), kiểm tra treo, crash, lỗi bộ nhớ, kết quả bất thường.
  \item \textbf{Sửa lỗi thực thi}: tiếp tục sửa tối thiểu, ghi lại nguyên nhân
        gốc và cách khắc phục.
  \item \textbf{Kiểm tra tính đúng}: mở ảnh, so sánh trực quan với ảnh tham
        chiếu; với mã MPI, kiểm tra cả trường hợp một tiến trình và nhiều tiến
        trình.
\end{enumerate}

Mỗi lỗi được ghi vào nhật ký với các trường: mô hình, biến thể, loại lỗi (biên
dịch/thực thi/đúng đắn), thông điệp lỗi, nguyên nhân gốc, cách sửa, số dòng thay
đổi và mức độ khó. Cách ghi này cho phép phân loại lỗi ở Phần~\ref{sec:ketqua}
và tách \emph{triệu chứng} khỏi \emph{nguyên nhân}.

Một điểm phương pháp luận cần nêu rõ: sau khi sửa, chương trình được đo hiệu
năng. Vì vậy kết quả mở rộng phản ánh chất lượng của “mã sau sửa”, không phải
chất lượng nguyên bản. Điều này hợp lý nếu mục tiêu là đánh giá năng suất thực
tế — người dùng luôn phải sửa trước khi dùng — nhưng cần tránh diễn giải quá mức
thành “LLM sinh mã chạy đúng ngay”.

\subsection{Độ đo}

\subsubsection{Tỉ lệ thành công}

Với mỗi mô hình, tỉ lệ thành công của từng tiêu chí được tính trên toàn bộ năm
chương trình:
\begin{equation}
  R_{\text{crit}} = \frac{\#\{\text{chương trình đạt tiêu chí crit}\}}{5}\times100\%.
\end{equation}
Cách tính này cho phép so sánh trực tiếp giữa các mô hình, nhưng che khuất sự
khác biệt giữa các biến thể. Vì vậy báo cáo trình bày thêm bảng lỗi chi tiết
theo từng biến thể.

\subsubsection{Kích thước và chi phí mã}

Số dòng mã (LOC) được đo bằng công cụ \code{cloc}, bao gồm cả chú thích, trước
và sau khi sửa. Chi phí phát triển được ước lượng bằng mô hình COCOMO \cite{boehm1984} thông qua
công cụ \code{scc}. Công thức cơ bản của COCOMO:
\begin{equation}
  \text{Effort} = a \cdot (\text{KLOC})^{b},
\end{equation}
với $a,b$ phụ thuộc loại dự án. COCOMO được thiết kế cho phần mềm tuần tự, không
mô hình hóa chi phí đồng bộ, truyền thông hay gỡ lỗi song song. Vì vậy con số
COCOMO chỉ nên đọc như một thước đo tương đối giữa các chương trình cùng loại,
không phải dự báo tuyệt đối.

\subsubsection{Thông lượng và tăng tốc}

Thông lượng được đo bằng số cập nhật điểm ảnh mỗi giây
(\term{pixel updates per second}). Với ảnh $W\times H$ và $I_{\max}$ vòng lặp,
một cách hiểu thận trọng là đại lượng tỉ lệ nghịch với thời gian tính toán của
cùng một khối lượng công việc:
\begin{equation}
  \text{Throughput} \;\propto\; \frac{W\cdot H\cdot I_{\max}}{T_{\text{compute}}}.
\end{equation}
Khi so sánh các đường cong, cần chú ý cả \emph{độ dốc} (khả năng mở rộng) lẫn
\emph{độ cao tuyệt đối} (hiệu năng nền). Một đường phẳng nằm cao có thể là dấu
hiệu chương trình luôn dùng toàn bộ lõi, không phải dấu hiệu mở rộng tốt.

Tăng tốc và hiệu suất được tính theo Công thức~\eqref{eq:speedup}. Với thí
nghiệm phân tán, $P$ là số nút (1, 2, 4, 8, 12), mỗi nút một rank MPI và 20 lõi
OpenMP.

\subsection{Môi trường thực nghiệm}

\begin{table}[H]
\centering
\small
\caption{Cấu hình phần cứng và phần mềm của thực nghiệm.}
\label{tab:env}
\begin{tabularx}{\textwidth}{@{}p{4.2cm} X@{}}
\toprule
\textbf{Thành phần} & \textbf{Cấu hình} \\
\midrule
Cụm & rostam, Đại học Bang Louisiana; phân vùng medusa gồm 16 nút \\
CPU & Intel Xeon Gold 6148 @ 2,40 GHz, 20 lõi/nút \\
Bộ nhớ & 92 GB/nút \\
Trình biên dịch & GCC 11.5.0 \\
MPI & OpenMPI 5.0.7 \\
Backend cho C++17 & Intel TBB 2022.2 \\
Ảnh bộ nhớ chia sẻ & $1920\times1080$ điểm ảnh, 1000 vòng lặp \\
Ảnh phân tán & $12000\times80000$ điểm ảnh, 1000 vòng lặp \\
Cấu hình đo phân tán & 1, 2, 4, 8, 12 nút; 1 rank MPI/nút; 20 lõi/nút \\
Công cụ đo mã & \code{cloc}, \code{scc} \\
\bottomrule
\end{tabularx}
\end{table}

Việc dùng một cấu hình duy nhất giúp so sánh công bằng nhưng giới hạn khả năng
ngoại suy. Kết quả có thể khác trên Clang, trên OpenMP thật (không qua TBB), trên
GPU hoặc trên kiến trúc many-core. Phần~\ref{sec:thaoluan} thảo luận các mối đe
dọa này.

\subsection{Tính tái lập và kiểm soát chất lượng}

Mã sinh ra và mã tham chiếu được công bố trên GitHub và Zenodo \cite{diehl2026}.
Việc công bố cả mã lỗi, bản sửa và ảnh đầu ra cho phép người đọc tái lập từng
bước. Trong quá trình đánh giá, nhóm tác giả giữ nguyên prompt, thứ tự chạy và
cách đo; chỉ thay đổi những dòng cần thiết để chương trình hoạt động. Mọi thay
đổi được ghi lại, nhờ đó có thể phân biệt lỗi của mô hình với can thiệp của con
người.

Một hạn chế cần nêu là mỗi cấu hình chỉ được đo một lần. Bài báo không báo cáo
số lần lặp, độ lệch chuẩn hay khoảng tin cậy. Với các phép đo có biến thiên nhỏ
như thông lượng trung bình trên 20 lõi, sai số có thể chấp nhận được; nhưng với
các điểm gần biên (ví dụ bước nhảy 2$\to$4 nút), cần thận trọng khi diễn giải.
Đây là lý do báo cáo tập trung vào \emph{xu hướng} và \emph{nguyên nhân} hơn là
so sánh chênh lệch nhỏ giữa các đường cong.
